\documentclass{article}
\usepackage{graphicx} % Required for inserting images
\usepackage{color, xcolor}
\usepackage{amsmath}
\usepackage{amssymb}
\usepackage{booktabs}
\usepackage{hyperref}
\usepackage{geometry}
\geometry{margin=1in}

\title{Wavelet--Surrogate Non-Stationarity Detector\\
\large Gatekeeper v1-F: Multi-Moment Analysis + Multichannel Consistency Gate}
\author{Mario De Los Santos, Luis Enrique Sucar, Felipe Orihuela-Espina}
\date{March 2026}

\begin{document}

\maketitle
\tableofcontents
\newpage

% ==============================================================================
\section{Goal}
Given a multi-channel time series
\[
X(t)=\{x_1(t),\dots,x_N(t)\}, \quad t=1,\dots,T,
\]
the objective is to partition the timeline into intervals where the signal is
\emph{approximately stationary} so that a causal snapshot can be computed per interval.
Gatekeeper flags a set of change points
\[
\tau_1,\dots,\tau_K
\]
and corresponding small event neighborhoods
\[
[t_{\text{start}}^k,\, t_{\text{end}}^k],\quad k=1,\dots,K,
\]
that contain statistically significant non-stationarity transitions.

For each detected event, the detector reports:
(i) the event type (onset or offset),
(ii) contributing channels ranked by energy score,
(iii) contributing scales/frequencies (optionally),
(iv) a peak derivative magnitude $\mathrm{peak\_dE}$,
(v) a local area severity, and
(vi) per-channel and per-moment energy breakdowns.

% ==============================================================================
\section{Core Hypothesis: Surrogate-Based Null Testing}
We avoid ad-hoc change-point rules sensitive to noise drift and instead use a surrogate null.

\begin{itemize}
  \item \textbf{Null hypothesis} $H_0$: the observed signal (or derived feature) is indistinguishable from a stationary linear stochastic process with the \emph{same power spectrum} as a baseline segment.
  \item \textbf{Alternative hypothesis} $H_1$: the signal contains \emph{time-localized structures} (amplitude/frequency modulations or distributional changes) that cannot be generated by such a stationary process.
\end{itemize}

Key idea: thresholds are \emph{data-driven} via baseline surrogates, so if the noise floor increases globally the reference increases too, reducing false positives.

The detector is two-sided: it detects both \emph{onset} events (energy increase relative to baseline, $dE > \epsilon_+$) and \emph{offset} events (energy decrease, $-dE > \epsilon_-$), corresponding to transitions into and out of non-stationary regimes.

% ==============================================================================
\section{Preprocessing: Robust Artifact Rejection}
Before calibration and feature extraction, each channel $x_{ch}(t)$ is cleaned with a robust outlier rule based on Median Absolute Deviation (MAD).

\paragraph{Algorithm.}
\begin{enumerate}
  \item Compute the median $\tilde{x}$ and MAD of the finite samples:
  \[
    \mathrm{MAD} = \mathrm{Median}\!\left(\left|x(t) - \tilde{x}\right|\right).
  \]
  \item Scale to be comparable to a standard deviation (consistent estimator under Gaussianity):
  \[
    \hat{\sigma}_{\mathrm{MAD}} = 1.4826\cdot\mathrm{MAD}.
  \]
  \item Flag samples exceeding the threshold:
  \[
    \mathrm{bad}(t) = \mathbf{1}\!\left[\left|x(t)-\tilde{x}\right| / \hat{\sigma}_{\mathrm{MAD}} > \delta_{\mathrm{MAD}}\right]
    \quad\text{or}\quad \mathbf{1}[\neg\mathrm{finite}(x(t))],
  \]
  with default $\delta_{\mathrm{MAD}} = 5.0$.
  \item Replace flagged samples by linear interpolation over the nearest good neighbors.
\end{enumerate}

This reduces baseline-spectrum inflation and prevents isolated outliers from masking subtle events.

% ==============================================================================
\section{Feature Construction: Rolling Statistical Moments (v1-E)}
v1-E detects nonstationarity not only in raw signals but also in the \emph{distributional moments} of rolling windows. This extends sensitivity beyond spectral changes in the raw signal to changes in mean, spread, shape, and tail behavior.

\paragraph{Moment time series.}
For each channel $x_{ch}(t)$ (after artifact rejection) and rolling window length $w_m$, define:
\[
m_{ch}^{(1)}(t)=\mathrm{mean}_{w_m}(t),\quad
m_{ch}^{(2)}(t)=\mathrm{var}_{w_m}(t),\quad
m_{ch}^{(3)}(t)=\mathrm{skew}_{w_m}(t),\quad
m_{ch}^{(4)}(t)=\mathrm{kurt}_{w_m}(t),
\]
where skewness is the standardized 3rd central moment and kurtosis is excess kurtosis (Fisher's definition).

Efficient computation via vectorized rolling windows:
\[
m_{ch}^{(1)}(t) = \frac{1}{w_m}\sum_{\tau} x_{ch}(\tau),\quad
m_{ch}^{(2)}(t) = \mathbb{E}[x^2]_{w_m} - \left(\mathbb{E}[x]_{w_m}\right)^2.
\]
Skewness and kurtosis are computed using a strided sliding-window view for efficiency:
\[
z(\tau) = \frac{x_{ch}(\tau) - \hat{\mu}_{w_m}}{\hat{\sigma}_{w_m}},\quad
m^{(3)}(t) = \frac{1}{w_m}\sum z^3,\quad
m^{(4)}(t) = \frac{1}{w_m}\sum z^4 - 3.
\]

\paragraph{Causal vs.\ centered windows.}
A causal (backward-looking) rolling window can be used to avoid ``early'' detection around change points:
at time $t$, uses samples $[t - w_m + 1, t]$.
Centered windows (default: $[t - w_m/2, t + w_m/2]$) can be used when delay is acceptable.
The implementation uses reflect-padding to handle boundaries.
Default: $w_m = 50$ samples, causal mode.

\paragraph{Moment weighting (factorial-inverse).}
Moment contributions are aggregated with weights
\[
w_m = \frac{1}{m!},\quad m \in \{1, 2, 3, 4\},
\]
so lower-order moments dominate (more stable) while higher-order moments still contribute.
Concretely: $w_1 = 1.0$, $w_2 = 0.5$, $w_3 \approx 0.167$, $w_4 \approx 0.042$.

% ==============================================================================
\section{Signal Transform: CWT with Complex Morlet}
For a univariate series $y(t)$ (either a raw channel or a rolling moment series), compute the
Continuous Wavelet Transform (CWT) using a complex Morlet mother wavelet:
\[
W(a,t)=\frac{1}{\sqrt{a}}\int y(\tau)\,\psi^*\!\left(\frac{t-\tau}{a}\right)\,d\tau,
\]
with scalogram energy
\[
\mathrm{SG}(a,t)=|W(a,t)|^2,
\]
over a scale set $\mathcal{A}$ (log-spaced, with minimum scale to suppress high-frequency noise).

\paragraph{Morlet wavelet.}
The complex Morlet wavelet at scale $s$ and center frequency $\omega_0$ is:
\[
\psi(t; s, \omega_0) = \frac{1}{\sqrt{s}}\,e^{j\omega_0 t/s}\,e^{-t^2/(2s^2)},
\]
implemented as FFT-based convolution (\texttt{scipy.signal.fftconvolve}) with reflect-padding of
$3\cdot a_{\max}$ samples to avoid boundary artifacts. Default $\omega_0 = 6.0$.

\paragraph{Scale set.}
Scales are log-spaced between $a_{\min}$ and $a_{\max}$:
\[
\mathcal{A} = \mathrm{logspace}(\log_2 a_{\min},\, \log_2 a_{\max},\, n_{\mathrm{scales}}),\quad
a_{\max} = \min\!\left(64,\, \max\!\left(\lfloor T/20 \rfloor, 8\right)\right),
\]
with $n_{\mathrm{scales}} = \min(20, \max(5, \lfloor \log_2(a_{\max}/a_{\min})\cdot 4\rfloor))$ and
default $a_{\min} = 4.0$.

\paragraph{Approximate frequency mapping (Morlet).}
For interpretability, scales map to pseudo-frequencies as
\[
f \approx \frac{\omega_0 \, f_s}{2\pi\, a},
\]
where $f_s$ is the sampling rate.

% ==============================================================================
\section{Phase A: Calibration (Building the Noise Reference)}
\textbf{Input:}
a baseline index set $\mathcal{B}\subset\{1,\dots,T\}$, number of surrogates $N_{\text{surr}}$,
and significance level $\alpha$.

% -----------------------------------------------------------------------
\subsection{Fourier Surrogates}
For each baseline feature series $y_{\text{base}}(t)$, build surrogates preserving the FFT magnitude
while randomizing the phase:
\[
y_{\text{surr}}^{(i)}(t) = \mathcal{F}^{-1}\!\Big\{|Y(f)|\,e^{j\phi^{(i)}(f)}\Big\},
\]
where $\phi^{(i)}(f)\sim \mathrm{Uniform}(0,2\pi)$ with conjugate symmetry enforced
($\phi[0] = 0$ for DC; $\phi[n/2] = 0$ for Nyquist if $n$ is even).

% -----------------------------------------------------------------------
\subsection{Per-Scale Scalogram Thresholds (Time-Quantile Method, v1-E)}
Max-over-time thresholds are too conservative (the surrogate maximum inflates with $T$, pushing
almost all observed SG below threshold, causing $E^{-}$ to dominate). Instead, v1-E uses a
\emph{time-quantile} at level $q = 0.95$.

Let $q_{\mathrm{SG}} = 0.95$. For each channel $ch$, moment $m$, and scale $a$, and each surrogate $i$:
\[
Q_{ch,m}^{(i)}(a) = \mathrm{Quantile}_{t\in\mathcal{B}}\!\Big(\mathrm{SG}_{ch,m}^{(i)}(a,t),\; q_{\mathrm{SG}}\Big).
\]
The per-scale calibration threshold is then:
\[
T_{ch,m}(a) = \mathrm{Quantile}_{i=1,\dots,N_{\mathrm{surr}}}\!\Big(Q_{ch,m}^{(i)}(a),\; 1-\alpha\Big).
\]
This yields a per-channel, per-moment, per-scale reference that balances sensitivity and specificity.

% -----------------------------------------------------------------------
\subsection{Transition Thresholds from Surrogate Peak Heights (v1-E)}
Detection operates on a \emph{differential} signal (derivative spikes at transitions), so thresholds
are calibrated on \emph{peak heights}, not global maxima or integrals (which inflate with $T$).

For each of $N_\epsilon = \min(N_{\mathrm{surr}}, 50)$ surrogate realizations:
\begin{enumerate}
  \item Build the weighted multi-moment signed energy $E_{\mathrm{signed}}^{(i)}(t)$ (see Section~\ref{sec:detection}).
  \item Compute the smoothed derivative $dE^{(i)}(t)$.
  \item Restrict to the interior (excluding the edge margin $l_e = \max(10, \lfloor a_{\max}\rfloor)$).
  \item Extract positive and negative peaks using \texttt{find\_peaks} with distance $= 20$:
  \[
    \mathcal{P}_+ = \bigcup_i \!\left\{\text{peak heights of } dE^{(i)}(t)\right\},\quad
    \mathcal{P}_- = \bigcup_i \!\left\{\text{peak heights of } -dE^{(i)}(t)\right\}.
  \]
\end{enumerate}
Then set:
\[
\epsilon_+ = \max\!\Big(\mathrm{Quantile}(\mathcal{P}_+,\; 1-\alpha),\; 0.1\Big),\qquad
\epsilon_- = \max\!\Big(\mathrm{Quantile}(\mathcal{P}_-,\; 1-\alpha),\; 0.1\Big).
\]
The floor of $0.1$ prevents spurious triggers from degenerate calibrations. If no peaks are found
in surrogates, the $99.9$th percentile of $|dE|$ is used as a fallback.

% ==============================================================================
\section{Phase B: Detection (Signed Wavelet Deviations and Energies)}
\label{sec:detection}

% -----------------------------------------------------------------------
\subsection{Signed Log-Ratio Deviation}
For each channel $ch$, moment $m$, compute the observed scalogram $\mathrm{SG}_{ch,m}(a,t)$ and
define a signed log-ratio deviation:
\[
Z_{ch,m}(a,t) = \log\!\left(\frac{\mathrm{SG}_{ch,m}(a,t)+\varepsilon}{T_{ch,m}(a)+\varepsilon}\right),
\quad \varepsilon = 10^{-10}.
\]
\begin{itemize}
  \item $Z_{ch,m}(a,t) > 0$: more energy than baseline threshold (potential onset).
  \item $Z_{ch,m}(a,t) < 0$: less energy than baseline threshold (potential offset).
\end{itemize}

Split into upward and downward components to prevent cancellation across scales:
\[
Z^{+}_{ch,m}(a,t) = \max(0, Z_{ch,m}(a,t)),\qquad
Z^{-}_{ch,m}(a,t) = \max(0, -Z_{ch,m}(a,t)).
\]

% -----------------------------------------------------------------------
\subsection{K-of-Scales Gate (Robustness)}
To suppress isolated artifacts at single scales, we optionally require at least $k$ scales to be
``significant'' at time $t$. Operationally:
\[
Z_{ch,m}(\cdot,t) \leftarrow 0 \quad \text{if fewer than } k \text{ scales satisfy } |Z_{ch,m}(a,t)| > \eta,
\]
with default $\eta = 0.5$ (in log-ratio units, corresponding to $\approx 1.65\times$ above or
$0.61\times$ below the threshold) and default $k = k_{\min} = 3$.

% -----------------------------------------------------------------------
\subsection{Per-Channel Moment Energies}
Collapse across scales:
\[
E^{+}_{ch,m}(t) = \sum_{a\in\mathcal{A}} Z^{+}_{ch,m}(a,t),\qquad
E^{-}_{ch,m}(t) = \sum_{a\in\mathcal{A}} Z^{-}_{ch,m}(a,t).
\]

% -----------------------------------------------------------------------
\subsection{Moment-Weighted Aggregation}
Aggregate energies across moments with factorial-inverse weights:
\[
E^{+}_{ch}(t) = \sum_{m} w_m\, E^{+}_{ch,m}(t),\qquad
E^{-}_{ch}(t) = \sum_{m} w_m\, E^{-}_{ch,m}(t).
\]

% -----------------------------------------------------------------------
\subsection{Channel Aggregation}
Aggregate across channels to obtain global energies, using one of three modes:
\[
E^{+}(t) =
\begin{cases}
\displaystyle\sum_{ch=1}^N E^{+}_{ch}(t) & (\text{\texttt{sum\_all}}), \\[6pt]
\displaystyle\max_{ch} E^{+}_{ch}(t) & (\text{\texttt{max}}), \\[6pt]
\displaystyle\sum_{ch\in\mathrm{TopK}} E^{+}_{ch}(t) & (\text{\texttt{topK}}),
\end{cases}
\]
and similarly for $E^{-}(t)$. Default: \texttt{sum\_all}.

The signed energy is:
\[
E_{\mathrm{signed}}(t) = E^{+}(t) - E^{-}(t).
\]

% ==============================================================================
\section{Energy Derivative and Forgetting Severity}

% -----------------------------------------------------------------------
\subsection{Differential Change Score}
To detect \emph{transitions} rather than sustained regimes, detection operates on the derivative
of a smoothed signed energy:
\[
\widetilde{E}(t) = \mathrm{MA}_{w_s}\!\left(E_{\mathrm{signed}}(t)\right),\qquad
dE(t) = \widetilde{E}(t) - \widetilde{E}(t-1),
\]
where $\mathrm{MA}_{w_s}$ denotes a causal moving average of window $w_s$ (default $w_s = 15$).
Within a stable regime, $E_{\mathrm{signed}}(t)$ may be elevated but $dE(t)\approx 0$, avoiding
repeated triggers.

% -----------------------------------------------------------------------
\subsection{Forgetting Severity Trace (Leaky Integrator)}
For visualization and a bounded notion of ``recent change mass'', compute
\[
\mathrm{sev}(t) = \lambda\,\mathrm{sev}(t-1) + |dE(t)|,\qquad \lambda\in(0,1),
\]
instead of an unbounded cumulative sum. Default $\lambda = 0.98$.
An alternative rolling-window mode ($\mathrm{sev}(t) = \sum_{\tau=t-w_r+1}^{t}|dE(\tau)|$) is
also available.

% ==============================================================================
\section{Peak-Based Event Extraction (v1-E)}
Events are extracted as peaks in $dE(t)$ (onsets) and peaks in $-dE(t)$ (offsets).

% -----------------------------------------------------------------------
\subsection{Candidate Peaks}
Using \texttt{scipy.signal.find\_peaks}, extract:
\[
\mathcal{C}_+ = \{t : dE(t) \text{ is a local maximum and } dE(t)\ge\epsilon_+\},\qquad
\mathcal{C}_- = \{t : -dE(t) \text{ is a local maximum and } -dE(t)\ge\epsilon_-\}.
\]
Parameters enforced:
\begin{itemize}
  \item A \textbf{refractory period} (minimum peak distance $d_{\min}$, default 200 samples) to avoid
        clustered duplicates.
  \item A \textbf{prominence} constraint: $\mathrm{prom}_{\pm} = 0.3\cdot\epsilon_{\pm}$, for
        robustness against slow drift.
  \item An \textbf{edge margin} $l_e = \max(10, \lfloor a_{\max}\rfloor)$ to exclude CWT boundary
        artifacts.
\end{itemize}

% -----------------------------------------------------------------------
\subsection{Step Validation (Sustained Level-Shift Filter)}
Derivative spikes can originate from transient artifacts. v1-E keeps a candidate $\tau$ only if it
corresponds to a sustained local step in $E_{\mathrm{signed}}$.

\paragraph{Local step filter.}
Let $\mathrm{pre}(\tau) = E_{\mathrm{signed}}[\tau - g - w_p\,:\,\tau - g]$ and
$\mathrm{post}(\tau) = E_{\mathrm{signed}}[\tau + g\,:\,\tau + g + w_p]$
where $g = 10$ (gap to exclude the transition itself) and $w_p = 160$ (window). Define
\[
\Delta(\tau) = \mathrm{mean}(\mathrm{post}(\tau)) - \mathrm{mean}(\mathrm{pre}(\tau)),
\qquad
\sigma_{\mathrm{pre}}(\tau) = \mathrm{std}(\mathrm{pre}(\tau)).
\]
Then:
\[
\text{onset: keep if } \Delta(\tau) \ge \kappa\,\sigma_{\mathrm{pre}}(\tau),\qquad
\text{offset: keep if } \Delta(\tau) \le -\kappa\,\sigma_{\mathrm{pre}}(\tau),
\]
with default $\kappa = 2.5$. This rejects short transients that do not shift the energy level.

\paragraph{Sequential step filter (alternative).}
For settings where multiple regime changes can produce cascading detections, a sequential variant
also tracks the \emph{current regime level} $L$, initialized to $\mathrm{median}(E_{\mathrm{signed}}[\mathcal{B}])$.
A candidate is accepted only if \emph{both} conditions hold:
\begin{align*}
\text{onset:} &\quad \mathrm{mean}(\mathrm{post}) - \mathrm{mean}(\mathrm{pre}) \ge \kappa\,\hat{\sigma}_{\mathrm{pre}} \quad \text{AND} \quad \mathrm{mean}(\mathrm{post}) - L \ge \kappa\,\hat{\sigma}_{\mathrm{pre}}, \\
\text{offset:} &\quad \mathrm{mean}(\mathrm{pre}) - \mathrm{mean}(\mathrm{post}) \ge \kappa\,\hat{\sigma}_{\mathrm{pre}} \quad \text{AND} \quad L - \mathrm{mean}(\mathrm{post}) \ge \kappa\,\hat{\sigma}_{\mathrm{pre}},
\end{align*}
where $\hat{\sigma}_{\mathrm{pre}}$ is the MAD-based robust sigma of the pre-window.
Upon acceptance, $L \leftarrow \mathrm{mean}(\mathrm{post})$, preventing multiple detections
within the same energy regime.

% -----------------------------------------------------------------------
\subsection{K-of-Channels Consistency Gate (v1-F)}
\label{sec:channel-gate}
A true structural regime change in a causal system should produce coordinated energy shifts across
multiple causally-connected channels (parent, child, and descendants are all affected by an edge
change). In contrast, a false positive caused by a single-channel noise artifact will typically
show a step in only one or two channels.

v1-F exploits this by adding a \emph{K-of-channels consistency gate}, the cross-channel analog of the
K-of-scales gate applied earlier per channel.

\paragraph{Per-channel step z-score.}
At each candidate peak $\tau$ (surviving step validation), compute for each channel $ch$:
\[
\Delta_{ch}(\tau) = \mathrm{median}\!\left(E^{\mathrm{signed}}_{ch}[\tau+g:\tau+g+w_p]\right)
 - \mathrm{median}\!\left(E^{\mathrm{signed}}_{ch}[\tau-g-w_p:\tau-g]\right),
\]
where $E^{\mathrm{signed}}_{ch}(t) = E^{+}_{ch}(t) - E^{-}_{ch}(t)$, with the same gap $g$ and
window $w_p$ as in step validation. Normalize by per-channel pre-window variability:
\[
z_{ch}(\tau) = \frac{\Delta_{ch}(\tau)}{\hat{\sigma}_{ch}(\tau)},
\qquad
\hat{\sigma}_{ch}(\tau) = 1.4826\cdot\mathrm{MAD}\!\left(E^{\mathrm{signed}}_{ch}[\tau-g-w_p:\tau-g]\right).
\]

\paragraph{Active channel count.}
Define the active indicator:
\[
A_{ch}(\tau) =
\begin{cases}
\mathbf{1}[z_{ch}(\tau) > \delta_{ch}] & \text{(onset)}, \\
\mathbf{1}[z_{ch}(\tau) < -\delta_{ch}] & \text{(offset)},
\end{cases}
\]
with default $\delta_{ch} = 1.5$ (MAD-scaled). The active count is $n_{\mathrm{active}}(\tau) = \sum_{ch} A_{ch}(\tau)$.

\paragraph{Gate rule.}
Keep candidate $\tau$ only if:
\[
n_{\mathrm{active}}(\tau) \ge K_{\mathrm{ch,min}},
\]
with default $K_{\mathrm{ch,min}} = 2$. For small $N$, the threshold is clamped:
$K_{\mathrm{ch,min}} \leftarrow \min(K_{\mathrm{ch,min}},\, \max(1, \lfloor N/2\rfloor))$.
The gate is auto-disabled for $N = 1$.

\paragraph{Concentration ratio (diagnostic).}
For post-hoc analysis, each accepted event stores:
\[
R(\tau) = \frac{\max_{ch}|\Delta_{ch}(\tau)|}{\sum_{ch}|\Delta_{ch}(\tau)| + \varepsilon}.
\]
$R \approx 1$ indicates the change is concentrated in a single channel (suspicious);
$R \approx 1/N$ indicates a distributed, multi-channel change (more likely genuine).

% -----------------------------------------------------------------------
\subsection{Signal-to-Noise Ratio (SNR) Filter}
As a final gating step, events are kept only if:
\[
\mathrm{SNR}(\tau) = \frac{\mathrm{peak\_dE}(\tau)}{\hat{\sigma}_{dE}^{(\mathcal{B})}} \ge \rho_{\min},
\]
where $\hat{\sigma}_{dE}^{(\mathcal{B})}$ is the standard deviation of $dE$ over baseline interior
samples (excluding edge margins), and $\rho_{\min} = 2.0$ by default.

% -----------------------------------------------------------------------
\subsection{Event Interval}
Each accepted change point $\tau$ is reported with a small neighborhood:
\[
[t_{\text{start}}, t_{\text{end}}] = [\tau - w,\; \tau + w],
\]
where $w = \max(3, w_s)$ (tied to the smoothing window).

% ==============================================================================
\section{Attribution and Event Scores}
For each event $k$ around $\tau_k$:

\paragraph{Channel attribution.}
For onset events, rank channels by accumulated $E^{+}_{ch}(t)$ over the event interval;
for offset events by $E^{-}_{ch}(t)$:
\[
S_{ch}^{(+)} = \sum_{t\in[t_s, t_e]} E^{+}_{ch}(t),\qquad
S_{ch}^{(-)} = \sum_{t\in[t_s, t_e]} E^{-}_{ch}(t).
\]
Channels are sorted in descending order to provide a contribution ranking.

\paragraph{Severity scores.}
Two severity metrics are reported per event:
\[
\mathrm{peak\_dE} = \max_{t\in[t_s, t_e]} |dE(t)|,\qquad
\mathrm{area} = \sum_{t\in[t_s, t_e]} |dE(t)|.
\]
The forgetting severity trace $\mathrm{sev}(t)$ provides a global, time-continuous visualization
of ``recent change mass''; it is \emph{not} used for event extraction.

\paragraph{Per-moment breakdown (v1-E).}
The signed energy per moment $E_{m,\mathrm{signed}}(t) = E^{+}_m(t) - E^{-}_m(t)$ is stored in
$\mathrm{E\_by\_moment}$, enabling post-hoc analysis of which statistical moment drives each event:
\[
\delta_m(\tau) = \mathrm{mean}_{t\in[\tau, \tau+w_p]}E_{m,\mathrm{signed}}(t)
 - \mathrm{mean}_{t\in[\tau-w_p, \tau]}E_{m,\mathrm{signed}}(t).
\]

% ==============================================================================
\section{Piecewise Recalibration (Optional)}
For long recordings with multiple regime changes, thresholds calibrated on the initial baseline
may become suboptimal after the signal has settled into a new regime. Piecewise recalibration
iteratively:
\begin{enumerate}
  \item Runs the full detector (initial detection).
  \item Splits the signal at detected change points into segments.
  \item For each segment, selects a ``quiet'' subset as the samples with the lowest $|dE|$
        (lowest $q_{\mathrm{quiet}} = 0.70$ quantile), as a local baseline.
  \item Re-runs detection on each segment with local calibration, translating segment-local
        event indices back to global time.
  \item Repeats for $n_{\mathrm{iter}} = 2$ iterations (until change points stabilize).
\end{enumerate}
With v1-E's differential detection, recalibration is less critical than in level-based detectors
because $dE \approx 0$ inside sustained regimes. However, it remains useful for very long
recordings ($T > 5000$) with deviation jumps between regimes.

% ==============================================================================
\section{Detectability of Structural Changes}
A change in causal structure does not necessarily imply a strong change in univariate marginal
properties. In linear SEMs with i.i.d.\ noise, modifying edges can primarily alter cross-variable
dependence (e.g., covariance structure) while leaving marginal spectra/variances nearly unchanged.
Therefore, a detector based on per-channel wavelet energy (even multi-moment) may fail to highlight
some structural edits.

\paragraph{Detectability assumption.}
Structural regime shifts are detectable by Gatekeeper v1-E only when at least one channel exhibits
a change in mean, variance, skewness, kurtosis, and/or band-power dynamics; not purely in
cross-channel dependence structure.

% ==============================================================================
\section{Evaluation Protocol}

\subsection{Tolerance-Based Matching}
A detected change point $\hat{\tau}$ is a \emph{true positive} if there exists a ground-truth
change point $\tau^*$ within a tolerance window:
\[
|\hat{\tau} - \tau^*| \le \delta_{\mathrm{tol}}.
\]
Default $\delta_{\mathrm{tol}} = 100$ samples. Matching is performed greedily by minimum distance
(each ground-truth point matched to at most one detection).

\subsection{Metrics}
\[
\mathrm{Precision} = \frac{TP}{TP+FP},\quad
\mathrm{Recall} = \frac{TP}{TP+FN},\quad
F_1 = \frac{2\,\mathrm{Precision}\cdot\mathrm{Recall}}{\mathrm{Precision}+\mathrm{Recall}},
\]
\[
\mathrm{Accuracy} = \frac{TP}{TP+FP+FN},\quad
\mathrm{MTE} = \frac{1}{TP}\sum_{(gt, \hat\tau)\in\text{matches}} |gt - \hat\tau|.
\]
MTE (Mean Timing Error) measures localization accuracy among matched pairs.

\subsection{Post-Processing: Transition Simplification}
Nearby detections that likely correspond to the same physical transition are merged:
\[
\text{if } |\hat{\tau}_{j+1} - \hat{\tau}_j| < d_{\min},\quad
\hat{\tau}_j \leftarrow \left\lfloor\frac{\hat{\tau}_j + \hat{\tau}_{j+1}}{2}\right\rfloor.
\]
Default $d_{\min} = 50$ samples.

% ==============================================================================
\section{Default Hyperparameters and Implementation Notes}

\subsection{Core Parameters}

\begin{center}
\begin{tabular}{lll}
\toprule
\textbf{Parameter} & \textbf{Default} & \textbf{Description} \\
\midrule
$a_{\min}$ (\texttt{min\_scale}) & 4.0 & Minimum CWT scale (suppresses HF noise) \\
$\omega_0$ (\texttt{omega0}) & 6.0 & Morlet center frequency \\
$N_{\mathrm{surr}}$ (\texttt{n\_surrogates}) & 100 & Number of Fourier surrogates \\
$\alpha$ (\texttt{alpha}) & 0.01 & Significance level (controls threshold strictness) \\
$k_{\min}$ (\texttt{k\_scales\_min}) & 3 & Minimum significant scales (K-of-scales gate) \\
$w_s$ (\texttt{smooth\_window}) & 15 & Moving-average window for dE smoothing \\
$d_{\min}$ (\texttt{refractory\_period}) & 200 & Minimum distance between detected peaks \\
$\rho_{\min}$ (\texttt{min\_snr}) & 2.0 & Minimum SNR for accepting events \\
$\lambda$ (\texttt{severity\_decay}) & 0.98 & Leaky integrator decay \\
$w_m$ (\texttt{moment\_window}) & 50 & Rolling window for moment computation \\
$\delta_{\mathrm{MAD}}$ (\texttt{mad\_threshold}) & 5.0 & MAD outlier rejection threshold \\
$\kappa$ (\texttt{delta\_k}) & 2.5 & Step-validation threshold multiplier \\
$g$ (\texttt{gap}) & 10 & Gap between peak and pre/post windows \\
$w_p$ (\texttt{pre\_win, post\_win}) & 160 & Pre/post window for step validation \\
\midrule
\multicolumn{3}{l}{\textit{v1-F: Multichannel consistency gate}} \\
$K_{\mathrm{ch,min}}$ (\texttt{k\_channels\_min}) & 2 & Min channels with significant step \\
$\delta_{ch}$ (\texttt{delta\_ch\_k}) & 1.5 & Per-channel z-score threshold (MAD-scaled) \\
\bottomrule
\end{tabular}
\end{center}

\subsection{Experiment Configurations}
Three experiment configurations are provided in the evaluation runner:

\begin{center}
\begin{tabular}{llll}
\toprule
\textbf{Config} & \textbf{2-Regime} & \textbf{Multi-Regime} & \textbf{Realistic} \\
\midrule
$\alpha$ & 0.01 & 0.02 & 0.40 \\
$k_{\min}$ & 3 & 2 & 1 \\
$N_{\mathrm{surr}}$ & 75 & 75 & 100 \\
$w_s$ & 15 & 15 & 12 \\
$\rho_{\min}$ & 2.0 & 1.5 & 0.3 \\
$d_{\min}$ & 200 & $\min(150, n/5)$ & $\min(150, n_{\min}/4)$ \\
Baseline & $20\%$ of $T$ & $60\%$ of regime 1 & $50\%$ of regime 1 \\
\bottomrule
\end{tabular}
\end{center}

The realistic scenario uses a permissive $\alpha = 0.40$ and $\rho_{\min} = 0.3$ to maximize
recall across 5 regimes with variable lengths (600--800 samples each) and moderate structural
changes (25--35\%).

% ==============================================================================
\section{Complete Pipeline Summary (v1-F)}

The full detection pipeline is as follows:

\begin{enumerate}
  \item \textbf{Artifact rejection}: clean each channel with median+MAD interpolation ($\delta_{\mathrm{MAD}} = 5$).
  \item \textbf{Rolling moments}: compute $m_{ch}^{(1)}\dots m_{ch}^{(4)}$ for each channel with causal window $w_m$.
  \item \textbf{Calibration}: for each moment $m$ and channel $ch$:
        \begin{enumerate}
          \item Compute $N_{\mathrm{surr}}$ Fourier surrogates of the moment series.
          \item For each surrogate, compute the CWT scalogram and its $q_{\mathrm{SG}} = 0.95$ time quantile.
          \item Set $T_{ch,m}(a) = \mathrm{Quantile}_{i}(Q^{(i)}_{ch,m}(a),\, 1-\alpha)$.
        \end{enumerate}
  \item \textbf{Calibration of $\epsilon_{\pm}$}: generate $N_\epsilon$ surrogates, compute
        weighted $dE^{(i)}(t)$, collect peak heights, set $\epsilon_{\pm} = \mathrm{Quantile}(1-\alpha)$.
  \item \textbf{Signed deviations}: compute $Z_{ch,m}(a,t) = \log(\mathrm{SG}/T_{ch,m})$; apply K-of-scales gate.
  \item \textbf{Energy}: sum over scales $\to$ per-channel moment energies; weight by $1/m!$; aggregate channels.
  \item \textbf{Derivative}: smooth $E_{\mathrm{signed}}$ with $\mathrm{MA}_{w_s}$, then differentiate.
  \item \textbf{Severity}: leaky integrator on $|dE|$ with decay $\lambda$.
  \item \textbf{Peak extraction}: \texttt{find\_peaks} on $\pm dE$ with height $\epsilon_{\pm}$, distance $d_{\min}$, prominence ratio 0.3.
  \item \textbf{Step validation}: filter by local $\mathrm{mean}(\mathrm{post}) - \mathrm{mean}(\mathrm{pre}) \ge \kappa\,\sigma_{\mathrm{pre}}$.
  \item \textbf{K-of-channels gate} (v1-F): for each surviving peak $\tau$, compute per-channel
        $z_{ch}(\tau) = \Delta_{ch}/\hat{\sigma}_{ch}$; keep only if $n_{\mathrm{active}} \ge K_{\mathrm{ch,min}}$
        (Section~\ref{sec:channel-gate}).
  \item \textbf{SNR filter}: keep only $\mathrm{peak\_dE}/\hat{\sigma}_{dE}^{(\mathcal{B})} \ge \rho_{\min}$.
  \item \textbf{Event report}: build \texttt{Event} with $\tau$, type, ranked channels, $\mathrm{peak\_dE}$, area,
        and multichannel diagnostics ($n_{\mathrm{active}}$, $R$, $z_{ch}$).
\end{enumerate}

% ==============================================================================
\section{v1-F vs v1-E: Summary of Changes}

\begin{itemize}
  \item \textbf{New: K-of-channels consistency gate} (Section~\ref{sec:channel-gate}). For each
        candidate peak, computes per-channel MAD-scaled z-scores and requires
        $n_{\mathrm{active}} \ge K_{\mathrm{ch,min}}$ channels to show a significant step.
        Rejects false positives caused by single-channel noise artifacts.
  \item \textbf{New: Per-event multichannel diagnostics.} Each accepted event stores
        $n_{\mathrm{active}}$, $f_{\mathrm{active}} = n_{\mathrm{active}}/N$,
        per-channel $z_{ch}$ scores, and concentration ratio $R$.
  \item \textbf{Pipeline extension}: the channel gate is inserted between step validation and
        SNR filtering (step 11 in the pipeline summary).
  \item \textbf{Graceful degradation}: gate auto-disabled for $N=1$; threshold clamped to
        $\max(1, \lfloor N/2\rfloor)$ for small $N$.
\end{itemize}

\subsection{Features preserved from v1-E}
\begin{itemize}
  \item Multi-moment analysis (mean/variance/skewness/kurtosis) with factorial-inverse weights $1/m!$.
  \item Time-quantile SG thresholds ($q = 0.95$), peak-height $\epsilon_{\pm}$ calibration.
  \item Signed log-ratio deviations, K-of-scales gate, derivative-based peak detection.
  \item Step validation (local and sequential variants), SNR filter.
  \item Leaky integrator severity, median+MAD artifact rejection.
\end{itemize}

\end{document}
